\documentclass[11pt]{article}

\usepackage[a4paper,margin=0.92in]{geometry}
\usepackage{amsmath,amssymb,amsthm,mathtools}
\usepackage{graphicx}
\usepackage{xcolor}
\usepackage[colorlinks=true,linkcolor=blue!55!black,citecolor=blue!55!black,urlcolor=blue!55!black]{hyperref}
\usepackage{microtype}

\newtheorem{theorem}{Theorem}
\newtheorem{lemma}{Lemma}
\newtheorem{remark}{Remark}
\newcommand{\R}{\mathbb R}
\newcommand{\Sone}{\mathbb S^1}
\newcommand{\Hess}{\operatorname{Hess}}
\newcommand{\Var}{\operatorname{Var}}
\newcommand{\vol}{\operatorname{Vol}}

\title{A Two-Dimensional Counterexample to Convexity\\
Along Real $2$-Interpolations}
\author{Automatic Math Discovery Run \texttt{fa\_banach\_001}\\Model: GPT5.6}
\date{26 August 2026}

\begin{document}
\maketitle

\begin{center}
\fcolorbox{red!55!black}{red!3}{\parbox{0.91\textwidth}{
\textbf{Status.} Candidate full negative answer, likely valid, under the
nontrivial smooth-norm interpretation intended by the source.  An explicit
$C^\omega_+$ unconditional planar body gives a negative second variation;
Cauchy--Kowalevski turns it into an actual local real $2$-interpolation family
with nonconvex logarithmic partition function.
}}
\end{center}

\begin{abstract}
Cordero-Erausquin and Klartag conjectured that
$\alpha(t)=-\log\int_{\R^n}e^{-F_t}$ is convex along real
$2$-interpolation families whose endpoints are even, convex, and
$2$-homogeneous.  We disprove the conjecture already in dimension two.  Let
$K$ have support function
$h(\theta)=1-\frac1{100}\cos(4\theta)$ and let the homogeneous initial speed
have boundary value $a(\theta)=\cos(2\theta)$.  Cone coordinates reduce the
required second-variation inequality to
\[
 \int_0^{2\pi}h^2(a')^2\,d\theta
 \geq4\int_0^{2\pi}a^2h(h+h'')\,d\theta.
\]
The exact left-minus-right deficit is
$8\pi\varepsilon(3+4\varepsilon)=-148\pi/625<0$ for
$\varepsilon=-1/100$.  The homogeneous interpolation equation is an analytic
noncharacteristic PDE in angular coordinates, so the Cauchy--Kowalevski
theorem produces a short analytic $2$-interpolation through this datum.
After affine time rescaling, its two endpoints are unconditional smooth norm
squares on $[0,1]$, while $\alpha''(1/2)<0$.
\end{abstract}

\section{Source conjecture}

The source is Cordero-Erausquin and Klartag~\cite{CEK}, arXiv:1109.3652,
published in \emph{Geometric Aspects of Functional Analysis}, Lecture Notes
in Mathematics 2050 (2012), 151--168.  Conjecture 10 on PDF page 11 states
that if the endpoints of a real $2$-interpolation are even, convex, and
$2$-homogeneous norm squares, then the associated function $\alpha$ is
convex.

More precisely, the source asks whether, for
$F_i(x)=\lambda_i\|x\|_{K_i}^2$ with centrally symmetric convex bodies
$K_i$, ``the function $\alpha$ associated with the $2$-interpolation family
is convex.''  The source notes that evenness, convexity, and homogeneity
propagate along the interpolation.

\begin{figure}[ht]
 \centering
 \includegraphics[width=0.98\textwidth]{figures/open_problem_crop.png}
 \caption{Conjecture 10 and its motivation in the source, PDF page 11.}
\end{figure}

For a sufficiently regular family $F(t,x)$, real $2$-interpolation means
\begin{equation}\label{eq:pde}
 \partial_{tt}F
 =\frac12(\Hess_xF)^{-1}\nabla_x\partial_tF\cdot
 \nabla_x\partial_tF,
\end{equation}
and
\begin{equation}\label{eq:alpha}
 \alpha(t)=-\log\int_{\R^2}e^{-F(t,x)}\,dx.
\end{equation}
Differentiation under the integral gives, with
$d\mu_F=e^{-F}dx/\int e^{-F}$,
\begin{equation}\label{eq:secondvariation}
 \alpha''(t)=\frac12\int
 (\Hess F)^{-1}\nabla u\cdot\nabla u\,d\mu_F
 -\Var_{\mu_F}(u),\qquad u=\partial_tF.
\end{equation}

\section{Idea of the proof}

For a smooth strictly convex body $K\subset\R^2$, put
$F=\|\cdot\|_K^2/2$.  Polar integration separates the Gaussian radial
variable from the cone measure on $\partial K$.  If a $2$-homogeneous speed
$u$ has boundary value $a$, then the strengthened Brascamp--Lieb inequality
needed in~\eqref{eq:secondvariation} becomes a Poincar\'e inequality on the
centro-affine circle with spectral constant $4=2n$.

The Euclidean circle has exactly this constant, with equality for quadratic
speeds.  But the constant is not stable under nonellipsoidal perturbation.
For the fourth Fourier perturbation
$h=1+\varepsilon\cos4\theta$, the quadratic-looking test
$a=\cos2\theta$ has Rayleigh deficit
$8\pi\varepsilon(3+4\varepsilon)$.  A small negative $\varepsilon$ preserves
strict convexity and makes the deficit negative.  This proves negative
curvature of $\alpha$ at one time.  Analytic local existence for the
interpolation PDE then converts the infinitesimal datum into a genuine family
and hence a counterexample with two honest endpoints.

\section{The explicit counterexample}

\begin{theorem}\label{thm:main}
There exist $C^\omega_+$ unconditional centrally symmetric convex bodies
$K_0,K_1\subset\R^2$ and a real $2$-interpolation family
\[
 F_s(x)=\frac12\|x\|_{K_s}^2,\qquad 0\leq s\leq1,
\]
with endpoints $F_0,F_1$, such that the function in~\eqref{eq:alpha} is not
convex.  More precisely, the family may be chosen analytic on
$(0,1)\times(\R^2\setminus\{0\})$ and satisfies
$\alpha''(1/2)<0$.
\end{theorem}

We split the proof into the geometric computation and the local-existence
step.

\begin{lemma}[Cone-coordinate second variation]\label{lem:cone}
Let $K\subset\R^2$ be $C^\infty_+$ and centrally symmetric.  Write its
support function as $h(\theta)$ and its inverse Gauss parametrization as
\[
 X(\theta)=h(\theta)n(\theta)+h'(\theta)t(\theta),
 \quad n=(\cos\theta,\sin\theta),\quad
 t=(-\sin\theta,\cos\theta).
\]
Put $q=h+h''$, $F_*=\|\cdot\|_K^2/2$, and define the $2$-homogeneous
function $u_*$ by
\[
 u_*(rX(\theta))=r^2a(\theta),\qquad r>0.
\]
If $\int_0^{2\pi}a hq\,d\theta=0$, then the second variation prescribed by
the interpolation equation at the Cauchy datum $(F_*,u_*)$ is
\begin{equation}\label{eq:deficit}
 \alpha''(0)=\frac{A-4D}{M},
 \quad
 A=\int_0^{2\pi}h^2(a')^2\,d\theta,
 \quad D=\int_0^{2\pi}a^2hq\,d\theta,
 \quad M=\int_0^{2\pi}hq\,d\theta.
\end{equation}
\end{lemma}

\begin{proof}
The Jacobian of $x=rX(\theta)$ is
\[
 dx=r\det(X,X')\,dr\,d\theta=rh q\,dr\,d\theta.
\]
Euler homogeneity and differentiation along $\partial K$ give the orthogonal
identities
\[
 \Hess F_*(X,X)=1,\qquad \Hess F_*(X,X')=0,\qquad
 \Hess F_*(X',X')=\frac qh.
\]
Since $du_*(X)=2a$ and $du_*(X')=a'$, it follows that
\begin{equation}\label{eq:pointenergy}
 (\Hess F_*)^{-1}\nabla u_*\cdot\nabla u_*
 =r^2\left(4a^2+\frac hq(a')^2\right).
\end{equation}
The elementary radial integrals are
\[
 \int_0^\infty r e^{-r^2/2}\,dr=1,\quad
 \int_0^\infty r^3e^{-r^2/2}\,dr=2,\quad
 \int_0^\infty r^5e^{-r^2/2}\,dr=8.
\]
Thus $\int e^{-F_*}=M$, the mean of $u_*$ is zero, and
\[
 \Var_{\mu_{F_*}}(u_*)=\frac{8D}{M},\qquad
 \int(\Hess F_*)^{-1}\nabla u_*\cdot\nabla u_*\,d\mu_{F_*}
 =\frac{2(4D+A)}M.
\]
Substitution in~\eqref{eq:secondvariation} proves~\eqref{eq:deficit}.
\end{proof}

\begin{proof}[Proof of Theorem~\ref{thm:main}]
Take
\begin{equation}\label{eq:choice}
 h(\theta)=1+\varepsilon\cos4\theta,\qquad
 a(\theta)=\cos2\theta,\qquad \varepsilon=-\frac1{100}.
\end{equation}
Then
\[
 h>0,\qquad q=h+h''=1-15\varepsilon\cos4\theta>0,
\]
so $h$ is the support function of a $C^\omega_+$ unconditional body $K$.
Also $\int ahq=0$ by Fourier orthogonality.  Direct integration gives
\begin{align}
 M&=\pi(2-15\varepsilon^2),\label{eq:M}\\
 A&=\pi(4-4\varepsilon+2\varepsilon^2),\label{eq:A}\\
 D&=\pi\left(1-7\varepsilon-\frac{15}{2}\varepsilon^2\right).
 \label{eq:D}
\end{align}
Consequently,
\begin{equation}\label{eq:negative}
 A-4D=8\pi\varepsilon(3+4\varepsilon)
 =-\frac{148\pi}{625}<0.
\end{equation}
Lemma~\ref{lem:cone} therefore gives a negative formal second variation.

It remains to realize the Cauchy datum by an interpolation family.  In
ordinary polar coordinates $x=r(\cos\phi,\sin\phi)$, write
$F(t,x)=r^2f(t,\phi)$ and $v=f_t$.  Equation~\eqref{eq:pde} becomes
\begin{equation}\label{eq:angularpde}
 f_{tt}=
 \frac{2(2f+f_{\phi\phi})v^2-2f_\phi vv_\phi+f(v_\phi)^2}
 {2f(2f+f_{\phi\phi})-(f_\phi)^2}.
\end{equation}
The body in~\eqref{eq:choice} and its Gauss parametrization are analytic.
Moreover, the derivative of the Euclidean polar angle of $X(\theta)$ is
$hq/|X|^2>0$, so normal angle and polar angle are related by an analytic
diffeomorphism.  Hence $F_*$ and $u_*$ correspond to analytic periodic
Cauchy data $f(0,\phi)$ and $f_t(0,\phi)$.  The denominator in
\eqref{eq:angularpde} is the determinant of $\Hess F_*$ in the polar
orthonormal frame and is strictly positive.  The analytic
Cauchy--Kowalevski theorem~\cite{John} therefore supplies a periodic analytic
solution for $|t|<\tau_0$.  Compactness of the circle and uniqueness on
overlapping angular charts make $\tau_0$ uniform.

Positivity and strict convexity persist after decreasing $\tau_0$.  The
reduced ansatz makes $2$-homogeneity automatic, while uniqueness in the
analytic Cauchy problem preserves the coordinate reflection symmetries of the
data.  Hence every $F_t$ is an even,
unconditional, convex $2$-homogeneous function and can be written
$F_t=\|\cdot\|_{K_t}^2/2$ for a $C^\omega_+$ body $K_t$.
Differentiation under the integral is justified uniformly by the quadratic
upper and lower bounds on $F_t$.  Equations~\eqref{eq:deficit} and
\eqref{eq:negative} yield $\alpha''(0)<0$.

Finally choose $0<\tau<\tau_0$ and set
$\widetilde F_s=F_{(2s-1)\tau}$ for $0\leq s\leq1$.  Affine
time changes preserve~\eqref{eq:pde}, and
$\widetilde\alpha''(1/2)=4\tau^2\alpha''(0)<0$.  Taking
$K_0=K_{-\tau}$ and $K_1=K_\tau$ completes the counterexample.
\end{proof}

\section{Interpretation of smoothness at the origin}

A globally $C^2$ function on $\R^n$ that is exactly $2$-homogeneous is a
quadratic form: its Hessian at the origin determines the function by scaling.
Thus a literal reading of the source's preliminary ``smooth function on
$\R^n$'' definition would reduce Conjecture 10 to ellipsoids, where $\alpha$
is affine.  The conjecture explicitly writes its data as arbitrary norm
squares, so its only nontrivial interpretation is the standard one used for
smooth norms: $C^\infty$ and strongly convex on $\R^n\setminus\{0\}$,
continuous and convex at the origin, with the PDE holding off the origin (and
hence almost everywhere).  Theorem~\ref{thm:main} lies exactly in this
intended class.

\section{Verification and novelty boundary}

The algebra in~\eqref{eq:M}--\eqref{eq:negative} is independently checked by
the included SymPy script.  The sign has a large exact margin: the deficit is
$-148\pi/625$, while $h\geq99/100$ and $q\geq17/20$.  No numerical
approximation is used in the proof.

A bounded search on 26 August 2026 covered the run indexes, the exact arXiv id
and title, the phrases ``real $2$-interpolation,'' ``Rochberg--Semmes
interpolation,'' and ``$2$-homogeneous Brascamp--Lieb,'' and citations visible
from the source.  It found no later paper stating a proof or counterexample to
Conjecture 10.  Recent centro-affine Poincar\'e work proves inequalities with
constant $n$ for unconditional bodies~\cite{HuIvaki}; the conjecture here
requires the stronger constant $2n$, and our planar calculation shows that
this stronger constant fails.  Independent prior observation remains
possible, so the novelty label is candidate rather than certified.

The main human-review point is the standard analytic local-existence step for
the periodic homogeneous equation~\eqref{eq:angularpde}.  Everything before
that step is an exact two-line Fourier computation, and the PDE is explicitly
in noncharacteristic analytic normal form.

\begin{thebibliography}{9}

\bibitem{CEK}
D. Cordero-Erausquin and B. Klartag,
\emph{Interpolations, convexity and geometric inequalities},
in \emph{Geometric Aspects of Functional Analysis}, Lecture Notes in Math.
2050, Springer, 2012, 151--168;
\href{https://arxiv.org/abs/1109.3652}{arXiv:1109.3652}.

\bibitem{John}
F. John,
\emph{Partial Differential Equations}, 4th ed., Applied Mathematical
Sciences 1, Springer, 1982.

\bibitem{HuIvaki}
Y. Hu and M. N. Ivaki,
\emph{Centro-affine Poincar\'e inequality: unconditional convex bodies},
\href{https://arxiv.org/abs/2607.20223}{arXiv:2607.20223}.

\end{thebibliography}

\end{document}
